\begin{lexlean}{IdsRB}
\useglossary{lexlean.std.nat@1.2.0}
\useglossary{spec.titles@1.0.0}
\importmodule{Rows}
\title{Rust backend conformance claims}
\begin{semanticmodule}
definition{
  name: rb01,
  parameters: [],
  result: Rows::IdRow,
  body: (record Rows::IdRow {id := "RB-01", level := "build", suite := "rust-backend", statement := "Every Rust rendering is built as a closed AST each of whose constructs carries the calculus element it realizes, every correspondence row and every AST node is exercised by a run package, and a construct whose row, width, or program does not justify its element is refused.", sections := "§17.16"}),
  executable: true
}


definition{
  name: rb02,
  parameters: [],
  result: Rows::IdRow,
  body: (record Rows::IdRow {id := "RB-02", level := "build", suite := "rust-backend", statement := "An exported name that is not a lowercase snake-case identifier, is a Rust keyword, imitates a generated name, or is declared by the runtime, a name exported twice, and an unavailable crate name each fail with LLB6005, and a crate binding one name twice in a function is refused.", sections := "§17.16, §26.3"}),
  executable: true
}


definition{
  name: rb03,
  parameters: [],
  result: Rows::IdRow,
  body: (record Rows::IdRow {id := "RB-03", level := "build", suite := "rust-backend", statement := "An export that copies a parameter whose type is not Copy fails with LLB6005, and a crate that moves a value twice or reads it after moving it is refused.", sections := "§17.16, §26.3"}),
  executable: true
}


definition{
  name: rb04,
  parameters: [],
  result: Rows::IdRow,
  body: (record Rows::IdRow {id := "RB-04", level := "build", suite := "rust-backend", statement := "A package whose boundary holds a function value at any depth, whose program computes a value of a type no value inhabits, or whose rust-core program needs the heap fails with LLB6005, and a rust-core crate naming any heap type, runtime function, or construct is refused.", sections := "§17.13, §17.16"}),
  executable: true
}


definition{
  name: rb05,
  parameters: [],
  result: Rows::IdRow,
  body: (record Rows::IdRow {id := "RB-05", level := "build", suite := "rust-backend", statement := "A function that can overflow returns R<T> and every call to it propagates, any other returns its value, an export whose declared errors differ from its function's fails with LLB6005, and a crate that drops, invents, or misreturns a failure is refused.", sections := "§17.16"}),
  executable: true
}


definition{
  name: rb06,
  parameters: [],
  result: Rows::IdRow,
  body: (record Rows::IdRow {id := "RB-06", level := "build", suite := "rust-backend", statement := "Every committed package builds offline under its declared gates, rustc warnings and Clippy's default lints denied with ten documented exceptions, the export of each one with an observable outcome, called from a separate crate, prints exactly the denotation's observable outcome, and so does every primitive instance on its boundary and seeded inputs; a planted lint and a planted semantic mutation are detected.", sections := "§17.16"}),
  executable: true
}


definition{
  name: rb07,
  parameters: [],
  result: Rows::IdRow,
  body: (record Rows::IdRow {id := "RB-07", level := "build", suite := "rust-backend", statement := "Packages are deterministic and content-addressed: renderings by two separate processes under different roots and environments are byte-identical to each other and to the committed package, every manifest and provenance validates against its schema, the provenance binds the SHA-256 of each file, the program identity, the runtime LexLean's semantics records, the sources, and the language-1.2 compiler-semantics ID, and every committed package's sources are the semantic ID of the verified build stating its program.", sections := "§17.16, §21.7"}),
  executable: true
}


definition{
  name: rb08,
  parameters: [],
  result: Rows::IdRow,
  body: (record Rows::IdRow {id := "RB-08", level := "build", suite := "rust-backend", statement := "A read of a binding that nothing after it reads is a move and a projection of a pair component that nothing after it reads moves the component out of the pair, every other read clones, and every clone a rendering keeps is one the ownership check requires; a list is a run of contiguous cells, and building or taking apart a cell is constant work on any list; a crate that reads or moves a binding after a component was moved out of it, or moves a component twice, is refused.", sections := "§17.16"}),
  executable: true
}


definition{
  name: rb09,
  parameters: [],
  result: Rows::IdRow,
  body: (record Rows::IdRow {id := "RB-09", level := "build", suite := "rust-backend", statement := "Every name the Rust runtime takes from Rust's library is a row of its library table, each a machine primitive with a stated meaning, stated in the closed AST over primitives, or a construct of the closed AST, and every primitive is cited by a row; a runtime that names anything else, a row it does not use, a row citing a name that is not a primitive, and a primitive no row cites are refused.", sections := "§17.16"}),
  executable: true
}


definition{
  name: rb10,
  parameters: [],
  result: Rows::IdRow,
  body: (record Rows::IdRow {id := "RB-10", level := "build", suite := "rust-backend", statement := "The scaffolding of a Rust package is data and not renderer text: the edition, whether the crate is published, the lint gates, and every allowed lint with its Clippy group and reason are read from the language-1.2 semantic module `RustPackage`, `#![no_std]` is the production target row's `allocation = false`, and `#![forbid(lint)]` restates each gate on rustc's lints at `forbid`; a renderer source holding a string of the module, a module that is not the shape the reader accepts, and a `Cargo.toml` or crate header that differs from the module are refused.", sections := "§17.16"}),
  executable: true
}


definition{
  name: rb11,
  parameters: [],
  result: Rows::IdRow,
  body: (record Rows::IdRow {id := "RB-11", level := "build", suite := "rust-backend", statement := "Every lint a package allows states its Clippy group, none is in the `correctness` group, the module states the same as a theorem over the same list, which the reader checks by its shape and the module's project type-checks (`verify` replays its proof, this row does not), and the table of SPEC.md §17.16 equals the module's rows; an allowed lint in `correctness`, a specified row the module lacks, a module row the specification lacks, and a module without the theorem are refused.", sections := "§17.16"}),
  executable: true
}


definition{
  name: rb12,
  parameters: [],
  result: Rows::IdRow,
  body: (record Rows::IdRow {id := "RB-12", level := "build", suite := "rust-backend", statement := "The language-1.2 compiler-semantics ID covers the package scaffold: changing a lint row, the edition, or a gate of a copy of the embedded tree changes it, writing the module in its data form or reformatting its notation does not, and the module is outside the identities of languages 1.0 and 1.1.", sections := "§17.16, §21.2"}),
  executable: true
}


definition{
  name: rb13,
  parameters: [],
  result: Rows::IdRow,
  body: (record Rows::IdRow {id := "RB-13", level := "build", suite := "rust-backend", statement := "A package states the effects of its realization: `package.json` and `provenance.json` record them, ascending and registered, the production registry says that exhausting the stack aborts the process, and an unregistered or unordered list, or one that states allocation on `rust-core`, is refused.", sections := "§17.13, §17.16"}),
  executable: true
}


definition{
  name: rb14,
  parameters: [],
  result: Rows::IdRow,
  body: (record Rows::IdRow {id := "RB-14", level := "build", suite := "rust-backend", statement := "A library template whose self calls are all tail calls is rendered as a loop over its accumulator that breaks with its result, so that the production rows of the primitives it realizes carry no `recursion`, while a template with a call that is not a tail call is rendered as a call and keeps the effect; a rendering that states a self call where a loop is possible, and a row whose `recursion` differs from its rendering, are refused.", sections := "§17.13, §17.14, §17.16"}),
  executable: true
}

definition{
  name: rb15,
  parameters: [],
  result: Rows::IdRow,
  body: (record Rows::IdRow {id := "RB-15", level := "build", suite := "rust-backend", statement := "The Rust AST mirrors the machine callee and the overflow literal of RustSyntax: a crate whose functions call each machine primitive at each width is checked, printed, and compiled by the pinned rustc with warnings denied, each call prints as the Rust operation its primitive names, and each function computes on the boundary values of its width exactly what an independent computation over a wider integer states; a machine call with another number of operands than its primitive takes, and an overflow value outside the payload of an `Err`, are refused as LLB6005, and a lowered crate that holds a machine call is refused by the correspondence between a rendering and its program as a construct with no row, because machine calls appear only in the runtime profile.", sections := "§17.14, §17.16"}),
  executable: true
}

definition{
  name: rb16,
  parameters: [],
  result: Rows::IdRow,
  body: (record Rows::IdRow {id := "RB-16", level := "build", suite := "rust-backend", statement := "The core runtime every Rust package carries is the text of a base and of the functions the language-1.2 module `RustRuntime` defines: the compiler project's module, shipped byte for byte at `language/rust-runtime-1.2/RustRuntime.lex.tex`, defines `coreCrate`, a crate of `RustSyntax` that the Rust crate decodes, never evaluates, and prints in place of the handwritten function of the same name, and each printed function agrees with that handwritten function in result and in counted work on the boundary and seeded values of its parameters when both are compiled by the pinned rustc with warnings denied, and the decoded crate is checked by the validator that checks every rendering (`validate` and the normative `validateCrate`), with the runtime's reading of its arithmetic rule, so that a function of the core runtime fails only with an `Err` of the overflow value whose error type is the node `RustType.overflow` and allocates nothing; a module that names a constructor `RustSyntax` does not declare, gives a constructor a wrong number of arguments, refers to a definition it lacks, defines a datum with parameters, or lacks `coreCrate`, a function that returns an `Err` of another value or another error type, a heap type in the core runtime, a shipped copy that differs from the compiler project's module, a defined function that no differential covers, and a function whose body is changed to compute another result are each refused.", sections := "§17.16"}),
  executable: true
}

definition{
  name: rb17,
  parameters: [],
  result: Rows::IdRow,
  body: (record Rows::IdRow {id := "RB-17", level := "build", suite := "rust-backend", statement := "The core runtime's support items are the definitions of a second closed module: the language-1.2 module `RuntimeSupport` of the compiler project, shipped byte for byte at `language/rust-runtime-1.2/RuntimeSupport.lex.tex`, defines `supportCrate`, a crate of `RuntimeSyntax`, whose items are of a closed set (a `use`, a unit struct with derives, a type alias with parameters, a static atomic counter, a trait of method signatures, an impl of a trait or of a type with bounded parameters, a generic function with bounded parameters), whose expressions are of a closed set (a variable, a literal, a path, a call of a path, a method call, a field of a pair, a borrow, an equality, a match), and whose types, paths, and methods name a declared item or one entry of a closed table of the standard library; the Rust crate decodes it, never evaluates it, checks it with a validator of its own, and prints it where the handwritten text stood, so that the core runtime holds no handwritten Rust text, and `equal`, `compare`, `tick`, and `work` agree with the handwritten ones in result and in counted work on the boundary and seeded values of the core profile, and `compare` on seeded keys that nest pairs to depth two and three, when both are compiled by the pinned rustc with warnings denied; an item or expression outside the set, a name that is neither declared nor in the table, an impl that lacks a method of its trait or states another signature, a bound on a trait that is not declared, a heap type in the core runtime, a shipped copy that differs from the compiler project's module, and a changed body are each refused.", sections := "§17.16"}),
  executable: true
}

definition{
  name: rb19,
  parameters: [],
  result: Rows::IdRow,
  body: (record Rows::IdRow {id := "RB-19", level := "build", suite := "rust-backend", statement := "The closed module `RuntimeSyntax` carries the constructs the `rust-std` runtime is written in, and no others: items (a tuple struct, a record struct, and an enum with unit, tuple, and record variants, each generic and with derives, beside those of `RB-17`), types (a 128-bit width, `usize`, `char`, `str`, a slice, and a mutable reference), literals (a character and a string), patterns (a mutable binding, a tuple, a tuple struct, a struct, and a literal), and expressions (a `let` with an optional type, an assignment, an `if` with an `else`, a `while`, a `while let`, a `loop`, a `for`, a `break`, a `return`, a tuple, a struct literal, a binary operation over twelve operators, a negation, a cast, an index, a range slice, a dereference, a mutable borrow, a method call with type arguments, a block, and a `match` arm with a guard); the Rust crate decodes, validates, and prints a crate of them, the language-1.2 module `RuntimeProbe` of the compiler project, which holds one use of each construct and is shipped byte for byte, so that its functions, compiled by the pinned rustc with warnings denied and run on seeded values, return what an independent computation states; the validator names its rule for each refusal (a name that is neither declared nor in the closed table of the standard library, a call or method with another number of arguments, a variable not in scope, a `break` outside a loop, a struct literal that names a field the struct lacks, a tuple of one, a pattern that binds a name twice, an assignment to a name that is not mutable) and refuses in the core profile every type, path, and method of the table (51 names, which SPEC.md §17.16 lists and the case compares with the validator's) that allocates, as the hidden allocation the validator refuses in every `rust-core` crate; a constructor `RuntimeSyntax` lacks, a shipped copy that differs from the compiler project's module, and each refusal above are refused.", sections := "§17.16"}),
  executable: true
}

definition{
  name: rb20,
  parameters: [],
  result: Rows::IdRow,
  body: (record Rows::IdRow {id := "RB-20", level := "build", suite := "rust-backend", statement := "The std runtime every `rust-std` package carries is the text of a handwritten base and of the functions the language-1.2 module `RustStd` defines: the compiler project's module, shipped byte for byte at `language/rust-runtime-1.2/RustStd.lex.tex`, defines `stdCrate`, a crate of `RustSyntax` with the std profile that the Rust crate decodes, never evaluates, checks with the validator that checks every rendering, and prints in place of the handwritten function of the same name, and each printed function agrees with that handwritten function in result and in counted work on lists of every length from 0 to 70 and of 1,000, 4,096 and 200,000 cells, and on byte strings and texts of the empty, the single, the multi-byte, the invalid, and the long, built fresh, after an uncons has made a unique buffer longer than its run, shared by a clone, and with cells in front, when both are compiled by the pinned rustc with warnings denied; the changes this makes to the shipped crate are named and each has its plant: a type parameter of a signature (`RustType.param`, printed as the generic list `<T: Clone>`, a call of a function without it fails to compile), the machine primitives `Machine.listAppend`, `listLength`, `listIndex`, `listSlice`, `bytesAppend`, `bytesLength`, `stringLength`, `bytesIndex`, `bytesSlice`, `utf8Encode`, `utf8Decode`, `bytesCompare`, `stringSplit`, and `stringJoin` (families 26 to 39, printed as calls of the methods of `List`, `Bytes`, and `Str` that hold the walks, for example `List::length(&a)` and `Bytes::slice(&a, start, count)`; a call with two operands of the one-operand length is refused by the validator), the decoder of the std crate (`stdCrate` of the std profile, a module without it or with another profile is refused), and the base's methods of `List`, `Bytes`, and `Str` that hold the walks, a miscount and an extra tick of which the differential finds; a defined function that no differential covers, a function of the oracle that the generated runtime lacks, a shipped copy that differs from the compiler project's module, and a function whose body is changed are each refused.", sections := "§17.16"}),
  executable: true
}


definition{
  name: rb22,
  parameters: [],
  result: Rows::IdRow,
  body: (record Rows::IdRow {id := "RB-22", level := "build", suite := "rust-backend", statement := "A fallible function rendered as a loop that is its value returns the operand of each `Err` that ends the loop, so a call of another fallible function there, looped or not, is in tail position and returns the callee's result without `?`: the module of issue 279 and a variant whose callee is itself a loop render in `rust-std`, export, and replay, their packages realize the programs the case lowers and agree with the calculus evaluator on every input of the case, an overflow included, and the same call bound before the loop's end without propagating, or left as the loop body's value with the end dropped, is refused because the fallible call does not propagate its overflow.", sections := "§17.13, §17.16"}),
  executable: true
}
\end{semanticmodule}
\end{lexlean}
